% Options for packages loaded elsewhere
\PassOptionsToPackage{unicode}{hyperref}
\PassOptionsToPackage{hyphens}{url}
\PassOptionsToPackage{dvipsnames,svgnames,x11names}{xcolor}
\documentclass[
  11pt,
  a4paper,
]{article}
\usepackage{xcolor}
\usepackage[margin=19mm,headheight=15pt]{geometry}
\usepackage{amsmath,amssymb}
\setcounter{secnumdepth}{-\maxdimen} % remove section numbering
\usepackage{iftex}
\ifPDFTeX
  \usepackage[T1]{fontenc}
  \usepackage[utf8]{inputenc}
  \usepackage{textcomp} % provide euro and other symbols
\else % if luatex or xetex
  \usepackage{unicode-math} % this also loads fontspec
  \defaultfontfeatures{Scale=MatchLowercase}
  \defaultfontfeatures[\rmfamily]{Ligatures=TeX,Scale=1}
\fi
\usepackage{lmodern}
\ifPDFTeX\else
  % xetex/luatex font selection
\fi
% Use upquote if available, for straight quotes in verbatim environments
\IfFileExists{upquote.sty}{\usepackage{upquote}}{}
\IfFileExists{microtype.sty}{% use microtype if available
  \usepackage[]{microtype}
  \UseMicrotypeSet[protrusion]{basicmath} % disable protrusion for tt fonts
}{}
\makeatletter
\@ifundefined{KOMAClassName}{% if non-KOMA class
  \IfFileExists{parskip.sty}{%
    \usepackage{parskip}
  }{% else
    \setlength{\parindent}{0pt}
    \setlength{\parskip}{6pt plus 2pt minus 1pt}}
}{% if KOMA class
  \KOMAoptions{parskip=half}}
\makeatother
\setlength{\emergencystretch}{3em} % prevent overfull lines
\providecommand{\tightlist}{%
  \setlength{\itemsep}{0pt}\setlength{\parskip}{0pt}}

\usepackage{fontspec}
\setmainfont{Liberation Serif}
\setsansfont{Liberation Sans}
\setmonofont{Liberation Mono}[Scale=0.86]
\usepackage{ragged2e,indentfirst,xurl}
\usepackage{longtable,booktabs,array,calc}
\usepackage{xcolor,fancyhdr,titlesec,enumitem}
\definecolor{ink}{HTML}{19354A}
\definecolor{accent}{HTML}{167D8D}
\titleformat{\section}{\Large\sffamily\bfseries\color{ink}}{}{0pt}{}
\titleformat{\subsection}{\normalsize\sffamily\bfseries\color{ink}}{}{0pt}{}
\titleformat{\subsubsection}{\normalsize\sffamily\bfseries\color{ink}}{}{0pt}{}
\titlespacing*{\section}{1.25cm}{0pt}{8pt}
\titlespacing*{\subsection}{1.25cm}{8pt}{4pt}
\titlespacing*{\subsubsection}{1.25cm}{6pt}{3pt}
\setlength{\parindent}{1.25cm}
\setlength{\parskip}{3pt}
\setlength{\emergencystretch}{3em}
\setlength{\tabcolsep}{4pt}
\renewcommand{\arraystretch}{1.08}
\setlist{nosep,leftmargin=1.5em,topsep=3pt}
\pagestyle{fancy}\fancyhf{}
\fancyhead[L]{\small\sffamily ICSI438 / SEMINAR 02}
\fancyhead[R]{\small\sffamily Nogoolin / M2}
\fancyfoot[R]{\small\thepage}
\renewcommand{\headrulewidth}{0.3pt}
\AtBeginDocument{\fontsize{10.5}{12.5}\selectfont\justifying\setlength{\parindent}{1.25cm}}
\usepackage{bookmark}
\IfFileExists{xurl.sty}{\usepackage{xurl}}{} % add URL line breaks if available
\urlstyle{same}
\hypersetup{
  colorlinks=true,
  linkcolor={Maroon},
  filecolor={Maroon},
  citecolor={Blue},
  urlcolor={blue!55!black},
  pdfcreator={LaTeX via pandoc}}

\author{}
\date{}

\newcounter{none}
\begin{document}

\section{SRS Scope Charter}\label{doc1-srs-scope-charter}

\textbf{Nogoolin · ICSI438 / M2 / UE-2 · Тэнгис · 2026-09-21 · v0.2,
review draft}

\subsection{Төслийг сольсон
шийдвэр}\label{doc1-ux442ux4e9ux441ux43bux438ux439ux433-ux441ux43eux43bux44cux441ux43eux43d-ux448ux438ux439ux434ux432ux44dux440}

Би Seminar 1-ийн ажлаа зуны дадлагын Document RAG Chatbot төсөл дээр
хийж, хүлээлгэн өгсөн. Week 2-оос өөрийн бодит бүтээгдэхүүн, portfolio
болох Nogoolin-ийг үндсэн кейсээр сонгов. Энэ төсөлд шаардлага,
архитектур, API, deployment-ийн баримтыг одоогийн кодтой нийцүүлэх
хэрэгцээ байгаа тул M4, M5, M8, M9-ийн ажлыг бодит төсөлдөө үргэлжлүүлэх
боломжтой. Seminar 1-ийг дахин хийхгүй; шилжилтийн persona холбоог
M2-ийн нэмэлтээр ил тод тэмдэглэнэ.

\textbf{Public repository:}
\href{https://github.com/Tengis01/Nogoolin}{github.com/Tengis01/Nogoolin}

\subsection{Зорилго ба
уншигч}\label{doc1-ux437ux43eux440ux438ux43bux433ux43e-ux431ux430-ux443ux43dux448ux438ux433ux447}

Каталогийн хэрэглэгч бүтээгдэхүүний тухай асуулга илгээх, хадгалсан
бүтээгдэхүүнээ дахин харах боломжийг тодорхойлно. Гол уншигч нь
хэрэгжүүлж, шалгах хөгжүүлэгч Тэнгис; хоёрдогч уншигч нь хичээлийн багш.
Phase 0 баримтыг эх материал болгон ашиглаж, технологийн сонголтыг
SDD-д, шалгаж болох үр дүнг SRS-д тусгана.

\subsection{Included --- энэ
хүрээнд}\label{doc1-included-ux44dux43dux44d-ux445ux4afux440ux44dux44dux43dux434}

\begin{itemize}
\tightlist
\item
  \textbf{Inquiry + wishlist:} guest болон нэвтэрсэн хэрэглэгчийн
  асуулга, хүлээн авсан баталгаа; нэвтэрсэн хэрэглэгчийн
  хадгалах/харах/хасах урсгал. Админ inbox нь асуулга очсоныг шалгах
  төгсгөлийн цэг.
\item
  \textbf{Өгөгдлийн тусгаарлалт:} хэрэглэгч өөрийн wishlist, inquiry
  history-г л уншиж/өөрчилнө; админ эрхийг тусдаа шалгана.
\item
  \textbf{Hero fallback:} 3D asset бэлэн бус үед орлуулагч, WebGL-гүй
  болон хөдөлгөөн багасгах тохиргоонд статик хувилбар; каталогт хүрэх
  боломж.
\item
  \textbf{Шалгах ба build хийх хүрээ:} API-ийн local Docker build,
  deploy-оос өмнөх тестийн зайлшгүй хаалт. Docker-ийн сонголт SDD-д,
  амжилт/алдааны шалгуур NFR-03-д байна.
\end{itemize}

\subsection{Excluded --- M2/M3-ийн энэ багцаас
хассан}\label{doc1-excluded-m2m3-ux438ux439ux43d-ux44dux43dux44d-ux431ux430ux433ux446ux430ux430ux441-ux445ux430ux441ux441ux430ux43d}

\textbf{Бүтээгдэхүүн бүрийн 360° viewer:} тус бүрд нь 3D asset бэлтгэх
ажил ганцаараа хөгжүүлэх MVP-д хэт их тул энгийн зургийн gallery
ашиглана (D-07). Энэ нь нүүрний 3D intro-г хассан шийдвэр биш.
Viewer-ийг post-MVP-д дахин авч үзнэ; энэ багцад хэрэгжүүлэх, тестлэх
үүрэг үүсгэхгүй.

\subsection{Postponed --- дараа
хийх}\label{doc1-postponed-ux434ux430ux440ux430ux430-ux445ux438ux439ux445}

\textbf{Phase 5:} захиалга, төлбөр/checkout, хүргэлтийн ажиллагаа; одоо
\nolinkurl{delivery_enabled=false} гэсэн шийдвэр хүчинтэй. \textbf{Phase
6:} cloud Supabase холбох, production deployment баталгаажуулах, App
Store/Play Store нийтлэх. Mobile inquiry/wishlist, бодит GLB/.riv asset
болон бүх каталог/admin шаардлагыг энэ таван шаардлага бүрэн хамрахгүй;
бүтээгдэхүүний backlog-д үлдээнэ.

\subsection{Хүрээ тогтоох үндэслэл ба дуусах
нөхцөл}\label{doc1-ux445ux4afux440ux44dux44d-ux442ux43eux433ux442ux43eux43eux445-ux4afux43dux434ux44dux441ux43bux44dux43b-ux431ux430-ux434ux443ux443ux441ux430ux445-ux43dux4e9ux445ux446ux4e9ux43b}

Chinchilla-ийн х. 53--55 дахь аргаас: \textbf{хэн юу хүлээж байна?} ---
хэрэглэгч баталгаа, хөгжүүлэгч давтаж шалгах нөхцөл; \textbf{ямар далд
таамаг байна?} --- нэвтэрсэн л бол бүх өгөгдөл харах эрхтэй биш, asset
болон cloud бэлэн гэж үзэхгүй; \textbf{эхлээд юуг баримтжуулах вэ?} ---
эхлэлээс үр дүн хүртэлх inquiry/wishlist урсгал. Эдгээр нь номын аргыг
төсөлдөө хэрэглэсэн асуултууд, үгчилсэн эшлэл биш.

M2-ийн бэлэн багц: SRS/SDD загвар, 5 шаардлага, 8 баганат traceability,
persona холбоо, review draft. Ярилцлага хийгдээгүй; peer review илгээх,
нийтлэх/submission алхмыг хэрэглэгчийн заавраар алгассан. W1 persona-ийн
шинэ төсөлд шилжүүлсэн холбоог reviewer-ээр баталгаажуулах шаардлагатай.

\textbf{Эх:} \href{../../phase-0/01-vision.md}{Vision},
\href{../../phase-0/10-roadmap.md}{Roadmap},
\href{../../../agent-context/DECISIONS.md}{D-07 / ADR-011},
\href{Software_Project_Documentation_Week_02_Handout.pdf}{W2 handout};
Chinchilla, бүлэг 5, хэвлэмэл х. 53--55 (локал номын PDF х. 74--76).


\end{document}
